\documentclass[11pt,letterpaper]{article}
\pagestyle{empty}
\usepackage{amsmath, amssymb, amsthm}
\usepackage{latexsym, epsfig, ulem, cancel, multicol, hyperref}
\usepackage{graphicx, tikz, subfigure,pgfplots}
\usepackage[margin=1in]{geometry}
\setlength{\parindent}{0pt}
\usepackage{multirow}
\usepackage{mathtools}
\usepackage{verbatim}
\usepackage{tikz}
\usepackage{pgfplots}
\setlength{\parskip}{1ex}

\newcommand{\T}[0]{\top}
\newcommand{\F}[0]{\bot}
\newcommand{\liminfty}[1]{\lim_{#1 \to \infty}}
\newcommand{\limzero}[1]{\lim_{#1 \to 0}}
\newcommand{\limto}[1]{\lim_{#1}}
\newcommand{\Z}{\mathbb{Z}}
\newcommand{\R}{\mathbb{R}}
\newcommand{\C}{\mathbb{C}}
\newcommand{\Q}{\mathbb{Q}}
\newcommand{\odd}[0]{\mathbb{Z} - 2\mathbb{Z}}
\newcommand{\lineint}[1]{\int_{#1}}
\newcommand{\pypx}[2]{\frac{\partial #1}{\partial #2}}
\newcommand{\divg}{\nabla \cdot}
\newcommand{\curl}{\nabla \times}
\newcommand{\dydx}[2]{\frac{\text{d} #1}{\text{d} #2}}
\newcommand{\sqbkt}[1]{\left[ #1 \right]}
\newcommand{\paren}[1]{\left( #1 \right)}
\newcommand{\tribkt}[1]{\left< #1 \right>}
\newcommand{\abso}[1]{\left|#1 \right|}
\newcommand{\zero}{\{0\}}
\newcommand{\then}{\rightarrow}
\newcommand{\nonneg}{\Z^+ \cup \{0\}}
\DeclarePairedDelimiter\ceil{\lceil}{\rceil}
\DeclarePairedDelimiter\floor{\lfloor}{\rfloor}
\newcommand{\union}[2]{\bigcup_{#1}^{#2}}
\newcommand{\inter}[2]{\bigcap_{#1}^{#2}}
\newcommand{\openclose}[1]{\left( #1 \right]}
\newcommand{\closeopen}[1]{\left[ #1 \right)}
\newcommand{\compo}[2]{#1 e^{i #2}}
\newcommand{\laplase}{\bigtriangleup}
\newcommand{\bra}[1]{\left< #1 \right|}
\newcommand{\ket}[1]{\left| #1 \right>}
\newcommand{\braket}[2]{\left< #1 \mid #2 \right>}
\newcommand{\ketbra}[2]{\left| #1 \right> \left< #2 \right|}
\newcommand{\ketpsit}{\ket{\psi(t)}}
\newcommand{\ketphit}{\ket{\phi(t)}}
\newcommand{\ham}{\mathbf{H}}
\newcommand{\unx}{\hat{\mathbf{x}}}
\newcommand{\uny}{\hat{\mathbf{y}}}
\newcommand{\uns}{\hat{\mathbf{s}}}
\newcommand{\unr}{\hat{\mathbf{r}}}
\newcommand{\untheta}{\hat{\boldsymbol\theta}}
\newcommand{\unphi}{\hat{\boldsymbol\phi}}
\newcommand{\dfl}{\text{d}}

\newcommand{\wsnumber}{1}
\newcommand{\wstopic}{Vectors}
\pgfplotsset{
    every linear axis/.append style={
       axis x line=center,
       axis y line=center,
       xlabel={$x$},
       ylabel={$y$}
    },
    every axis plot/.append style={thick,mark=none}
}
\tikzset{
    point/.style={circle,draw,fill,minimum width=0.3ex,inner sep=0pt,outer sep=0pt},
    every label/.append style={black}
}


\usepackage[margin=1in]{geometry}
\usepackage{amsmath, amssymb, amsthm, graphicx, hyperref}
\usepackage{enumerate}
\usepackage{fancyhdr}
\usepackage{multirow, multicol}
\usepackage{tikz}
\pagestyle{fancy}
\fancyhead[RO]{Dennis Li}
\fancyhead[LO]{Analytical Mechanics }
\usepackage{comment}
\newif\ifshow
\showfalse

\ifshow
  \newenvironment{solution}{\textbf{Solution.}}{}
\else
  \excludecomment{solution}
\fi

\renewcommand{\thefootnote}{\fnsymbol{footnote}}
\usepackage{comment}


\newtheorem*{remark}{Remark}


\begin{document}
\begin{center}
\ifshow
  \textbf{\Large Problem Set 01 Solution}\\
\else
  \textbf{\Large Problem Set 01}\\
\fi
Instructor \\ Prof. Gabe\\
\end{center}

\hrule

\vspace{0.2cm}

\section{Rock it}
\subsection*{The Rocket Equation}
\paragraph{Solution}
\begin{enumerate}[i.]
    \item Consider a bag around the rocket and the ejected mass $\text{d}m$, we can write the total momentum inside the bag as
    \[
    \mathbf{p}(t) = \paren{M+\text{d}m}\mathbf{v}(t)
    \]
    Where $\mathbf{v}$ is the speed of the rocket in the lab frame. After the ejection, we have
    \[
    \mathbf{p}(t+\text{d}t) = \text{d}m\paren{\mathbf{u}+\mathbf{v}+\text{d}\mathbf{v}} + M\paren{\mathbf{v+\text{d}\mathbf{v}}}
    =\mathbf{u} +\mathbf{v}\text{d}m + M\mathbf{v}+M\text{d}\mathbf{v}
    \]
    Now we can find $\dfl\mathbf{p}$ as
    \[
    \dfl\mathbf{p} = \mathbf{u}\text{d}m + M\text{d}\mathbf{v} = \mathbf{F}_{ext}\cdot\dfl t = 0
    \]
    This equates to 0 since there are no external force in this bag system. Now we have
    \[
    \mathbf{u}\dydx{m}{t} + M\dydx{\mathbf{v}}{t} = 0
    \]
    Or we can rewrite it as
    \[
    \mathbf{u}\dot{M} + M\dot{\mathbf{v}} = 0
    \]
    To solve it, we can write it as
    \[
    -M\dfl\mathbf{v} = \mathbf{u}\dfl m 
    \]
    We can do some moving things around
    \[
    \dfl\mathbf{v} = \frac{\mathbf{u}}{M}\dfl m
    \]
    Now we can integrate both side
    \[
    -\Delta\mathbf{v} = \int_{m_0}^{m_f}\frac{\mathbf{u}}{M}\dfl m = \mathbf{u}\ln\frac{m_f}{m_0}
    \]
    Now we move the negative sign around and let the natural log eats it and flip the fraction
    \[
    \Delta \mathbf{v} = \mathbf{u}\ln\frac{m_0}{m_f}
    \]
    If the ship is initially at rest, we simply have
    \[
    \mathbf{v}(t) = \mathbf{u}\ln\frac{m_0}{m(t)}
    \]

    \item This is a trivial case, the initial speed $v_0$ makes no difference in our result since it is a constant and we are only interested in the change. If you are interested in the final speed, simply subtract it from the $\Delta \mathbf{v}$, or add it to the $\mathbf{v}(t)$

    \item In the case of changing mass carried and fuel carried, we are only interested in the natural log term since we want to achieve the same $\Delta \mathbf{v}$ We know that $m_0 - m_f = m_{fuel}$, and if we increased the payload by 25\%, in other word
    \[
    M'_0 = 1.25M_0
    \]
    Then to achieve the same ratio in the natural log, the final mass has to also increase by 25\%. This in turn tells us we need 25\% more fuel. The same logic is applied to that of 50\%. This is a linear increase. 

    \item Consider the rocket initially at rest, we will only be interested in the relationship of change in velocity and change in position, and we know that in this case
    \[
    \mathbf{v}(t) = \Delta\mathbf{v}
    \]
    And we can simply express the relationship as
    \[
    \mathbf{v}(t) = \mathbf{u}\ln\frac{m_0}{M(t)}
    \]
    Then, we can integrate
    \[
    \mathbf{y}(t) = \int_{0}^{t} \mathbf{v}(\tau)\dfl \tau = \int_{0}^{t}\mathbf{u}\ln\frac{m_0}{M(\tau)}\dfl \tau = \int_{0}^{t}\mathbf{u}\ln\frac{m_0}{M(\tau)}\dfl\tau
    \]
    Without information about how mass is changing, or how the rocket is propelling, we will have to leave the result as integral
    \[
    \mathbf{y}(t) = \int_{0}^{t}\mathbf{u}\ln\frac{m_0}{M(\tau)}\dfl\tau
    \]
\end{enumerate}
\paragraph{T1 3.10}
\textbf{Question:} Consider a rocket (initial mass $m_0$) accelerating from rest in free space. At first, as it speeds up, its momentum $p$ increases, but as its mass $m$ decreases, $p$ eventually begins to decrease. For what value of $m$ is $p$ maximum?

\textbf{Solution}

We know that
\[
\mathbf{p}(t) = m(t)\mathbf{v}(t)
\]
Since we are initially at rest, we can use the relationship we established before and get
\[
\mathbf{p}(t) = m\mathbf{u}\ln\frac{m_0}{m}
\]
Where $m$ is some function of time. Now we differentiate both side with respect to mass, since that is what we care about.
\[
\dydx{\mathbf{p}}{m} = \dydx{}{m}\paren{ m\mathbf{u}\ln\frac{m_0}{m}}
\]
This is equal to
\[
\dydx{\mathbf{p}}{m}  = \mathbf{u}\ln\frac{m_0}{m} - m\mathbf{u}\dydx{}{m}\paren{\ln m_0 - \ln m} = \mathbf{u}\ln\frac{m_0}{m} - \mathbf{u}
\]
We set this equal to $0$ to find the extrema. 
\[
\mathbf{u}\ln\frac{m_0}{m} - \mathbf{u} = 0
\]
In other word
\[
\mathbf{u}\ln\frac{m_0}{m} = \mathbf{u}
\]
This means that
\[
\ln\frac{m_0}{m} = 1
\]
And this in turn gives us
\[
m = \frac{m_0}{e}
\]
This is where the maxima happens. 



\paragraph{MT5 9-57}
A rocket in outer space in a negligible gravitational field starts from rest and accelerates uniformly at $a$ until its final speed is $v$. The initial mass of the rocket is $m_0$. How much work does the rocket’s engine do?

\paragraph{Solution} We know that the velocity of the rocket as a function of time is
\[
\mathbf{v}(t) = \mathbf{u}\ln\frac{m_0}{m(t)}
\]
Since we are interested in the total work done, we can simply say that velocity is a function of mass
\[
\mathbf{v}(m) = \mathbf{u}\ln\frac{m_0}{m}
\]
We know that kinetic energy is 
\[
K = \frac{1}{2}m\mathbf{v}^2
\]
And if we take the differential, with chain rule in mind, we have
\[
\dfl K = m\mathbf{v}\dfl \mathbf{v} + \frac{1}{2}\mathbf{v}^2\dfl m
\]
We can use the expression of velocity to find $\dfl v$
\[
\dfl \mathbf{v} = -\frac{\mathbf{u}}{m}\dfl m
\]
Therefore the expression for $\dfl K$ is given by 
\[
\dfl K =\frac{1}{2}\mathbf{v}^2\dfl m -\mathbf{u}\mathbf{v}\dfl m 
\]
Substitute the expression for velocity in, we have
\[
\dfl K =\frac{1}{2}\paren{\mathbf{u}\ln\frac{m_0}{m}}^2\dfl m - \mathbf{u}^2\ln\frac{m_0}{m}\dfl m 
\]
We can further simplify this expression to
\[
\dfl K = \paren{\ln\frac{m_0}{m}}\paren{\frac{1}{2}\ln\frac{m_0}{m}-1}\mathbf{u}^2\dfl m
\]
We can now integrate to find the total work done
\[
\int_0^{K_f} \dfl K=\mathbf{u}^2\int_{m_0}^{m_f}  \frac{1}{2}\ln^2\frac{m_0}{m} - \ln\frac{m_0}{m} \,\dfl m
\]
Now we can evaluate this integral, we see that the LHS is work $W$, which is what we want, the RHS gives us a massive expression. The final result is
\[
W = \frac{1}{2}m_f\mathbf{u}^2\ln^2\frac{m_0}{m_f}
\]


\section{Do not knock it til you rocket}
Treat the following problems as variable mass problem and make it like a rocket. 
\subsection{2 carts, some sand}
Material is blown into cart A from cart B at a rate $b$ kilograms per second, as shown. The material leaves the chute vertically downward, so that it has the same horizontal velocity $u$ as cart B. At the moment of interest, cart A has mass $M$ and velocity $v$. Find $\frac{dv}{dt}$, the instantaneous acceleration of A.
\begin{center}
    \includegraphics[width=0.35\linewidth]{Pictures//PS05/KK2-4-13.png}
\end{center}

\paragraph{Solution} We will wrap a bag around the truck A with a differential amount of sand $\dfl m$ coming in. This way, we can define the momentum of A at time t as
\[
\mathbf{p}(t) = M\mathbf{v} + \mathbf{u}\dfl m
\]
Since the sand is carrying some velocity $\mathbf{u}$ as it enters. Now, when the sand land in truck A after some differential amount of time $\dfl t$, we have
\[
\mathbf{p}(t+\dfl t) = (M+\dfl m)\paren{\mathbf{v}+\dfl\mathbf{v}} 
\]
Now we can find $\dfl \mathbf{p}$ as
\[
\dfl \mathbf{p} = \mathbf{p}(t+\dfl t) - \mathbf{p}(t) = M\mathbf{v}+M\dfl\mathbf{v}+\mathbf{v}\dfl m - M\mathbf{v} + \mathbf{u}\dfl m
\]
This is simply
\[
\dfl \mathbf{p} = M\dfl \mathbf{v}+\mathbf{v}\dfl m -\mathbf{u}\dfl m
\]
This gives 
\[
\dfl \mathbf{p} = M\dfl \mathbf{v}+\mathbf{v}\dfl m -\mathbf{u}\dfl m = \mathbf{F}_{ext}\cdot \dfl t = 0
\]
We can divide both side by $\dfl t$, and obtain
\[
\dydx{\mathbf{p}}{t} = M\dydx{\mathbf{v}}{t}+\mathbf{v}\dydx{m}{t} -\mathbf{u}\dydx{m}{t}
\]
We see that $b = \dydx{m}{t}$, therefore we have
\[
\dydx{\mathbf{p}}{t} = M\dydx{\mathbf{v}}{t}+ b\paren{\mathbf{v - \mathbf{u}}} = 0
\]
We have obtained the folloing ODE
\[
M\dot{\mathbf{v}} + b\mathbf{v} = b\mathbf{u}
\]
We define $\tau = \frac{M}{b}$, and rewrite the ODE as
\[
\dot{\mathbf{v}} + \frac{1}{\tau}\mathbf{v} = \frac{1}{\tau}\mathbf{u}
\]
The solution is simply
\[
\mathbf{v}(t) = \mathbf{v}_0e^{-\frac{t}{\tau}} - \mathbf{u}
\]



\subsection{Growing raindrop KK2-4.24}
A raindrop of initial mass $M_0$ starts falling from rest under the influence of gravity. Assume that the drop gains mass from the cloud at a rate proportional to the product of its instantaneous mass and its instantaneous velocity:
\[
\frac{dM}{dt} = k M \mathbf{v},
\]
where $k$ is a constant. 

Show that the speed of the drop eventually becomes effectively constant, and give an expression for the terminal speed. Neglect air resistance.


\paragraph{Solution}
We know that
\[
\mathbf{p}(t) = M\mathbf{v} 
\]
Therefore we know that
\[
\dydx{\mathbf{p}}{t} = M\dydx{\mathbf{v}}{t}+\mathbf{v}\dydx{M}{t}
\]
Substitute what we know about the rate of change of mass, we have
\[
\dydx{\mathbf{p}}{t} = M\dydx{\mathbf{v}}{t}+\mathbf{v}\paren{kM\mathbf{v}} = M\dydx{\mathbf{v}}{t}+kM\mathbf{v}^2 
\]
We also know that
\[
\mathbf{F}_{ext} = M\mathbf{g} = \dydx{\mathbf{p}}{t} 
\]
Therefore we have obtained the following ODE
\[
M\dydx{\mathbf{v}}{t}+kM\mathbf{v}^2  = M\mathbf{g}
\]
Simplify, we have
\[
\dydx{\mathbf{v}}{t}+k\mathbf{v}^2  = \mathbf{g}
\]
We simply move stuff around to get
\[
\frac{1}{g-k\mathbf{v}^2}\dfl \mathbf{v} = \dfl t
\]
Integrating both side, we have
\[
\frac{1}{\sqrt{gk}}\tanh^-1\paren{\mathbf{v}\sqrt{\frac{k}{g}}} = t+ C
\]
We do some moving things around, and we have
\[
\mathbf{v}(t) = \sqrt{\frac{g}{k}}\tanh\paren{\sqrt{gk}\paren{t+C}}\, \hat{\mathbf{g}}
\]
We know the raindrop starts at rest, this is our initial condition. This gives
\[
\mathbf{v}(0) = 0 = \sqrt{\frac{g}{k}}\tanh\paren{\sqrt{gk}\paren{t+C}}\, \hat{\mathbf{g}}
\]
This shows that $C=0$, since there are no other solution to solve this. We have obtained the final solution
\[
\mathbf{v}(0) = 0 = \sqrt{\frac{g}{k}}\tanh\paren{t\sqrt{gk}}\, \hat{\mathbf{g}}
\]
For terminal velocity, we examine when $t\to\infty$, this gives us
\[
\mathbf{v}_{ter} = \sqrt{\frac{g}{k}}\, \hat{\mathbf{g}}
\]

\subsection{1 rope, 3 ways}
\subsection{Rope Hanging}
The rope that is hanging off the table has some length $x$, and the center of mass situates at $\frac{x}{2}$. We see that this part of the rope has mass $\lambda x$ where $\lambda$ is the linear mass density.

Now, we can express the momentum at time $t$, with the center of mass moving at some speed $v$
\[
\mathbf{p}(t) = \lambda x v
\]
And after some differential amount of time, we have
\[
\mathbf{p}(t+\dfl t) = \lambda x v + \lambda v \dfl x
\]
Now we can take the difference, and get
\[
\dfl \mathbf{p} = \lambda v \dfl x = \mathbf{F}_{ext} \cdot \dfl t
\]
Here, the external force is purely gravitational, so
\[
\dfl \mathbf{p} = \lambda v \dfl x = \lambda x g \cdot \dfl t
\]
We have arrived at
\[
v\dydx{x}{t}= gx
\]
We realize that this is just
\[
\dot{x}^2 - gx = 0
\]

\subsubsection{MT5 9-15}
We put a trash bag around the part of the rope that is dangling beneath the hole and a differential amount that is about to fall. Suppose the length of the rope that has fallen has length $x$, then we set up the initial momentum with some velocity $\mathbf{v}$
\[
\mathbf{p}(t) = \frac{M}{L_0}x\mathbf{v}
\]
And the momentum after a differential amount of time, here $L_0$ is the total length of the rope. The part $\dfl m$ would have joined the fall, and we can express it as
\[
\dfl m = \frac{M}{L_0}\dfl x
\]
Therefore the expression for momentum after the differential amount of time is
\[
\mathbf{p}(t+\dfl t) = \frac{M}{L_0}\paren{\mathbf{v} + \dfl \mathbf{v}}\paren{x + \dfl x} = \frac{M}{L_0}\paren{x\mathbf{v} + \mathbf{v}\dfl x + x\dfl \mathbf{v} + \dfl \mathbf{v}\dfl x}
\]
The second order differential simply vanishes, leaving only
\[
\mathbf{p}(t+\dfl t) = \frac{M}{L_0}\paren{\mathbf{v} + \dfl \mathbf{v}}\paren{x + \dfl x} = \frac{M}{L_0}\paren{x\mathbf{v} + \mathbf{v}\dfl x + x\dfl \mathbf{v}}
\]
We also realize that
\[
\mathbf{v} = \dydx{x}{t}
\]
Now we take the difference
\[
\dfl \mathbf{p} = \frac{M}{L_0}\paren{\mathbf{v}\dfl x  + x\dfl\mathbf{v}}
\]
We now include gravity, which makes it
\[
\frac{M}{L_0}\paren{\mathbf{v}\dfl x  + x\dfl\mathbf{v}} = \frac{M}{L_0}xg\cdot\dfl t
\]
This is essentially
\[
\mathbf{v}\dfl x  + x\dfl\mathbf{v} = xg\dfl t
\]
Now we substitute to express it purely in $x$, we have
\[
\mathbf{v}\dydx{x}{t} + x\dydx{\mathbf{v}}{t} = gx
\]
We This is essentially
\[
\dot{x}^2 + x\ddot{x} = gx
\]

We can inspect the dimension of the ODE, and realize that we need a function of $x(t) = \sqbkt{\text{m}}$, this means that the only thing can make sense to be the solution of this ODE is that
\[
x(t) = Agt^2
\]
Since this is the only combination that makes the correct dimension for our problem. Now we can shove this back into the ODE to get
\[
4A + 2A = 1
\]
This gives us the solution of the ODE
\[
x(t) = \frac{1}{6}gt^2 + x_0
\]
Where $x_0$ is our initial condition. 


\paragraph{MT5 9-19. Chain falling onto a table}
A chain such as the one in Example 9.2 (with the same parameters) of length $b$ and mass $\rho b$ is suspended from one end at a point that is a height $b$ above a table so that the free end barely touches the tabletop. At time $t = 0$, the fixed end of the chain is released. 

Find the force that the tabletop exerts on the chain after the original fixed end has fallen a distance $x$.

we put a trash bag around the table at the time moment where a differential amount of the chain is contacting the tabletop with mass $\dfl m$. We can write out the momentum of the bag assuming the center of mass is moving with some velocity $\mathbf{v}$, and the length of the rope that has contacted the table as $x$.

Stuff in the bag then have momentum at time $t$
\[
\mathbf{p}(t) = -\mathbf{v}\dfl m
\]
But after a differential amount of time, since its no longer moving, it will have
\[
\mathbf{p}(t) = 0
\]
We see that
\[
\dfl \mathbf{p} = \mathbf{v}\dfl m
\]
We also know that the falling speed of the chain can be expressed as
\[
\mathbf{v} = \dydx{x}{t}
\]
Therefore, we essentially have
\[
\dfl \mathbf{p} = \dydx{x}{t}\dfl m = \mathbf{F}_{ext}\cdot \dfl t
\]
The only force that is external here is the gravity on the fallen chain that is on the table and inside the bag, therefore 
\[
\dydx{x}{t}\dfl m =\lambda xg \cdot \dfl t
\]
We also notice that 
\[
\dfl m = \lambda \dfl x
\]
So we can further rewrite our equation as
\[
\lambda\dydx{x}{t}\dfl x =\lambda xg \cdot \dfl t
\]
And move the differential of time to the other side
\[
\lambda \dydx{x}{t}\dydx{x}{t} = \lambda xg
\]
This eventually trims down to
\[
\dot{x}^2 = xg
\]
And we have
\[
\dot{x}^2 - gx = 0
\]
We find that this is exactly the same as the first rope problem. And we can sort of eyeball the solution of this ODE. It will be something that involves $gt^2$, since we need to have unit of $[\text{m}]$, so lets take a guess
\[
x(t) = Agt^2
\]
We take the first derivative and square it to get
\[
\dot{x}^2 = 4A^2g^2t^2
\]
And the second term
\[
gx = Ag^2t^2
\]
We have
\[
4A^2 - A = 0 = A(4A-1)
\]
This means that
\[
A = \begin{cases}
    A = 0\\
    A = \frac{1}{4}
\end{cases}
\]
Here, we want the non-trivial solution, therefore the distance $x$ can be expressed as
\[
x(t) = \frac{1}{4}gt^2
\]
We can take the second derivative as
\[
\ddot{x} = a = \frac{g}{2}
\]
And the mass that piled up on the table is $\lambda x(t)$, therefore the normal force is simply
\[
F = \frac{g}{2}\cdot \lambda x(t) = \frac{\lambda}{8}g^2t^2
\]


\end{document}